\subsection{Сложные графики}
%%%%%%%%%%%%%%%%%%%%%%%%%%%%%%%%%%%%%%%%%%%%%%%%%%%%
%%%%%%%%%%%%%%%%%%%%%%%%%%%%%%%%%%%%%%%%%%%%%%%%%%%%
\z[{\begin{wrapfigure}[5]{r}{0.3\textwidth}%
		\vspace{-25pt}%
		\begin{flushright}%
 			\includegraphics[width=0.25\textwidth]{2.4.1}%
		\end{flushright}%
	\end{wrapfigure}%
	}]%
{На графике приведена зависимость скорости пешехода от времени его движения. С помощью этого графика постройте график зависимости скорости этого пешехода от расстояния, которое он прошёл. $v = 1.0\,\text{м/с}$.}
%
{}

%%%%%%%%%%%%%%%%%%%%%%%%%%%%%%%%%%%%%%%%%%%%%%%%%%%%
\z[{\begin{wrapfigure}[4]{r}{0.3\textwidth}%
		\vspace{-25pt}%
		\begin{flushright}%
 			\includegraphics[width=0.25\textwidth]{2.4.2}%
		\end{flushright}%
	\end{wrapfigure}%
	}]%
{На рисунке представлен график зависимости скорости тела от времени. Начальная координата тела $x(0) = 0$. Постройте графики зависимости ускорения и координаты тела, а также пройденного им пути от времени. Тело движется вдоль оси $Ox$.}
%
{}

%%%%%%%%%%%%%%%%%%%%%%%%%%%%%%%%%%%%%%%%%%%%%%%%%%%%
\z[{\begin{wrapfigure}[5]{r}{0.3\textwidth}%
		\vspace{-15pt}%
		\begin{flushright}%
 			\includegraphics[width=0.25\textwidth]{2.4.3}%
		\end{flushright}%
	\end{wrapfigure}%
	}]%
{На рисунке представлен график зависимости скорости тела от времени. Начальная координата тела $x(0) = 0$. Постройте графики зависимости ускорения и координаты тела, а также пройденного им пути от времени. Тело движется вдоль оси $Ox$.
\vspace{2em}
}
%
{}

%%%%%%%%%%%%%%%%%%%%%%%%%%%%%%%%%%%%%%%%%%%%%%%%%%%%
\z[{\begin{wrapfigure}[1]{r}{0.3\textwidth}%
		\vspace{-25pt}%
		\begin{flushright}%
 			\includegraphics[width=0.25\textwidth]{2.4.4}%
		\end{flushright}%
	\end{wrapfigure}%
	}]%
{По графику $a_x(t)$ (см. рисунок) постройте график $v_x(t)$, считая $v(0) = 0$.
\vspace{2em}
}
%
{}

%%%%%%%%%%%%%%%%%%%%%%%%%%%%%%%%%%%%%%%%%%%%%%%%%%%%
\z[{\begin{wrapfigure}[7]{r}{0.3\textwidth}%
		\vspace{-0pt}%
		\begin{flushright}%
 			\includegraphics[width=0.25\textwidth]{2.4.5}%
		\end{flushright}%
	\end{wrapfigure}%
	}]%
{На графике (рис.) представлена зависимость средней скорости машины от пройденного пути. Определите среднюю скорость машины на участке, где она разгонялась.}
%
{$120 \,\text{км/ч}.$}

%%%%%%%%%%%%%%%%%%%%%%%%%%%%%%%%%%%%%%%%%%%%%%%%%%%%
\z{Заяц убегает от Волка по прямой, двигаясь равномерно. В начальный момент времени расстояние между Зайцем и Волком равно $S = 36\,\text{м}$, а скорость Волка равна $v_0 = 14\,\text{м/с}$. Волк устаёт и через каждые $\Delta t = 10\,\text{с}$ (в моменты времени $\Delta t, 2\Delta t, 3\Delta t, \dots$, считая от начала движения) уменьшает свою скорость на $\Delta v = 1\,\text{м/с}$. С какой скоростью должен бежать Заяц, чтобы Волк его не поймал?}
%
{$11.8 \,\text{м/с}.$}
%%%%%%%%%%%%%%%%%%%%%%%%%%%%%%%%%%%%%%%%%%%%%%%%%%%%
\z[{\begin{wrapfigure}[8]{r}{0.4\textwidth}%
		\vspace{-15pt}%
		\begin{flushright}%
 			\includegraphics[width=0.4\textwidth]{2.4.7}%
		\end{flushright}%
	\end{wrapfigure}%
	}]%
{Автомобиль проехал по пятикилометровому участку дороги. Специальный прибор при этом записывал показания спидометра через каждые $10$ метров. В результате получилась зависимость скорости автомобиля $v$ от пройденного пути $x$, показанная на рисунке. Оцените, за какое время $t$ автомобиль проехал эти $5$ километров.}
%
{$(120 \pm 10) \,\text{с}.$}

